% ====================================================================
% Chapter 5: Evaluation Systems and Benchmarks (English expanded version)
% Target: 40KB+
% ====================================================================

\section{Evaluation Systems and Benchmarks}
\label{ch:evaluation}

\subsection{Evaluation Methodology}
\label{sec:eval-methodology}

\subsubsection{The Evaluation Challenge in Agent Self-Evolution}

The central tension in agent self-evolution research lies between a system's claim of improvement and our ability to verify that improvement is genuine, robust, and transferable. Evaluation is not merely the experimental section of a paper---it serves as the \emph{selection pressure} that shapes the direction of evolution and the \emph{safety gate} that prevents harmful mutations from propagating. If the evaluation function is misaligned, incomplete, or gameable, then self-evolution becomes self-deception.

This chapter addresses the evaluation problem from three perspectives. First, we propose a five-layer evaluation hierarchy that spans fully automated programmatic verification at the bottom to human-in-the-loop safety gates at the top. Second, we survey the major benchmarks used across the agent evolution literature, analyzing their strengths, limitations, and contamination risks. Third, we introduce agent-specific evaluation metrics that go beyond single-task accuracy to capture iterative improvement, transfer, cost efficiency, and safety. Throughout, we emphasize that evaluation infrastructure is itself evolving: as agents become more capable, evaluators must become more sophisticated, or risk being outpaced by the systems they are meant to assess.

The evaluation challenge is particularly acute for self-evolving systems because of four properties that distinguish them from conventional machine learning.

\textbf{Property 1: Moving targets.} An agent that modifies its own code, prompts, or memory may change its behavior in ways that the benchmark did not anticipate. The evaluated system at time $t+1$ is structurally different from the system at time $t$, making longitudinal comparison non-trivial. DGM~\citep{dgm2025} exemplifies this: each generation produces a qualitatively different agent architecture, and the SWE-Bench score at generation $n$ reflects a fundamentally different program than at generation $n-1$.

\textbf{Property 2: Evaluation as part of the optimization loop.} Agents can learn to exploit evaluator weaknesses, producing reward hacking, length bias, or score manipulation that inflates metrics without genuine capability gain. Self-Rewarding Language Models~\citep{selfrewarding2024} demonstrate both the promise and peril of this property: the model evaluates its own outputs, creating a feedback loop that initially produces genuine improvement but risks evaluator degradation over iterations.

\textbf{Property 3: Scope expansion.} The scope of evaluation must expand beyond task completion to include process quality, cost, safety, and transfer---dimensions that are harder to quantify than accuracy. An agent that solves 90\% of tasks but costs \$100 per run and occasionally deletes files is not necessarily superior to one that solves 80\% at \$0.10 per run with perfect safety.

\textbf{Property 4: Reproducibility threats.} Non-determinism in LLM generation, API version changes, and the interaction between stochastic sampling and evolving agent state threaten reproducibility. A self-evolving system tested with GPT-4-turbo-2024-04-09 may produce different results with GPT-4-turbo-2024-12-17, even if the agent code is identical. This version dependency is compounded by the evolutionary process itself: because each generation builds upon the results of previous generations, small differences in early generations can compound into qualitatively different evolutionary trajectories. Two runs of the same evolutionary system with different random seeds may discover entirely different solutions, making it difficult to distinguish robust improvements from lucky trajectories.

These four properties interact in subtle ways. Property 2 (evaluation in the loop) combined with Property 1 (moving targets) creates a situation where the evaluation landscape itself changes as the agent evolves, potentially invalidating conclusions drawn from early evaluation rounds. Property 3 (scope expansion) combined with Property 4 (reproducibility threats) means that multi-dimensional evaluation---already more expensive than single-metric evaluation---must be repeated across multiple random seeds and model versions, multiplying the cost of rigorous assessment.

These properties mean that a single benchmark score is never sufficient evidence for self-evolution. A credible evaluation must include multi-layer verification, cross-benchmark generalization, cost analysis, safety auditing, and ablation studies that isolate the contribution of the self-improvement mechanism from confounds such as prompt engineering, model selection, and data contamination.


\subsubsection{A Five-Layer Evaluation Hierarchy}

We propose organizing evaluation methods into five layers, from fully automated to human-supervised. Each layer adds reliability at the cost of increased effort and latency. A well-designed evaluation system uses all five layers in combination, with lower layers providing rapid feedback during development and higher layers serving as pre-deployment gates.

% ====================================================================
% Figure 5.1: Evaluation Benchmark Layered System
% ====================================================================
\begin{figure}[htbp]
\centering
\begin{tikzpicture}[
    layer/.style={rectangle, rounded corners=6pt, draw=#1!70!black, fill=#1!12,
        minimum width=12cm, minimum height=1.4cm, align=center, font=\small},
    bench/.style={font=\scriptsize, color=gray!60},
]
    % Layers from bottom to top
    \node[layer=cat1] (l1) at (0,0) {\textbf{L1: Programmatic Verification} --- Unit Tests, Type Checking, Linting, Fuzzing, Sandbox Execution};
    \node[layer=cat2] (l2) at (0,2) {\textbf{L2: Code \& Reasoning Benchmarks} --- SWE-Bench, HumanEval, GSM8K, MATH, MMLU, GPQA};
    \node[layer=cat3] (l3) at (0,4) {\textbf{L3: Environment Interaction Benchmarks} --- WebArena, ALFWorld, Voyager/Minecraft, AppWorld};
    \node[layer=cat4] (l4) at (0,6) {\textbf{L4: Process Metrics} --- Iterative Gain Curves, Archive Diversity, Transferability, Rollback Rate};
    \node[layer=cat5] (l5) at (0,8) {\textbf{L5: Safety \& Cost Gates} --- Permission Checks, Budget Limits, Regression Tests, Human Audit};

    % Side annotations
    \node[font=\scriptsize\bfseries, color=cat1!60!black, rotate=90] at (-7,0) {Automated};
    \node[font=\scriptsize\bfseries, color=cat5!60!black, rotate=90] at (-7,8) {Human};
    \draw[-{Stealth}, thick, dashed, color=gray!40] (-6.5,0) -- (-6.5,8);

\end{tikzpicture}
\caption{Five-layer evaluation hierarchy for agent self-evolution. From L1 (fully automated programmatic verification) to L5 (human-in-the-loop safety gates), evaluation cost and reliability increase at each layer.}
\label{fig:eval-layers}
\end{figure}

\textbf{L1: Programmatic Verification.} The most reliable evaluation signals come from deterministic programmatic checks---unit tests, type checking, static analysis, linting, fuzzing, regression testing, and sandbox execution. Systems like DGM~\citep{dgm2025}, SICA~\citep{sica2025}, and AlphaEvolve~\citep{alphaevolve2025} achieve self-evolution precisely because candidate code can be automatically executed and scored. FunSearch~\citep{funsearch2024} maintains 15 samplers and 140 evaluators (1:9 ratio) with a triple safety filter: execution success within 30-second timeout, no calls to ancestor functions, and return value type checking. AlphaEvolve produces discoveries such as a 48-multiplication algorithm for $4\times4$ complex matrices---the first improvement over Strassen's 49-multiplication result in 56 years---through pure programmatic evaluation.

\textbf{L2: Code and Reasoning Benchmarks.} Standardized benchmarks provide comparability across papers and systems. HumanEval~\citep{chen2021codex} (164 Python problems), MBPP~\citep{austin2021mbpp} (974 crowd-sourced problems), GSM8K~\citep{cobbe2021gsm8k} (8,500 math problems), and MATH~\citep{hendrycks2021math} (12,500 competition problems) are widely used but suffer from data contamination and limited task granularity.

\textbf{L3: Environment Interaction Benchmarks.} SWE-Bench~\citep{jimenez2024swebench} requires real open-source repository patching. WebArena~\citep{zhou2024webarena} creates realistic web-based tasks. ALFWorld~\citep{shridhar2021alfworld} tests embodied agents. Voyager~\citep{voyager2023} measures lifelong learning in Minecraft.

\textbf{L4: Process Metrics.} These evaluate \emph{how} improvement occurs: iterative improvement curves ($\Delta_t = V(z_t, \tau) - V(z_{t-1}, \tau)$), transfer and generalization, archive diversity, rollback rate, and cost efficiency.

\textbf{L5: Safety and Cost Gates.} Permission verification, cost accounting, regression testing, and human audit provide the final safety net for deployment decisions.


\subsection{Code Generation Benchmarks}
\label{sec:eval-code}

\subsubsection{HumanEval and HumanEval+}

HumanEval~\citep{chen2021codex} consists of 164 Python programming problems with unit tests, measuring functional correctness via pass@$k$ metrics. Each problem includes a function signature, docstring, body stub, and a set of unit tests. Despite its modest size, HumanEval has become the de facto standard for evaluating code generation capabilities.

Reflexion~\citep{reflexion2023} demonstrates one of the most striking results in self-improvement literature: 91\% pass@1 on HumanEval using GPT-3.5-turbo, surpassing GPT-4's 80.1\% zero-shot performance by 10.9 percentage points. The memory buffer---storing natural language descriptions of failures and corrections---acts as a ``verbal gradient'' guiding subsequent attempts. This result demonstrates that a smaller model with self-correction capabilities can outperform a larger model without them.

SelfEvolve~\citep{selfevolve2024} achieves 85.98\% pass@1 and 87.81\% pass@10 on HumanEval with a ChatGPT-based pipeline, improving from a 74.39\% baseline through combined self-evolution and self-refinement. Agent Symbolic Learning~\citep{agentsymbolic2024} achieves 73.2\% pass@1 with GPT-4 backbone, improving from the standard Agents baseline at 68.3\%.

HumanEval+ extends the original benchmark with additional test cases to address the well-known limitation that many solutions passing the original 164 unit tests may fail on edge cases not covered by the test suite. Absolute Zero's AZR-Coder-7B achieves SOTA at 7B parameter scale on HumanEval+ using zero external training data.

\textbf{Limitations.} HumanEval problems average approximately 10 lines of solution code with clearly specified inputs and outputs, while SWE-Bench tasks require understanding repository structure, navigating multiple files, and producing patches that integrate with existing code architecture. The granularity gap between these two evaluation settings is substantial: a method that shows 90\% improvement on HumanEval may show negligible improvement on SWE-Bench. Furthermore, HumanEval has been publicly available since 2021, making data contamination a significant concern for models trained on internet-scale corpora after this date.

\begin{table}[htbp]
\centering
\small
\caption{Code generation and software engineering benchmark results across self-improvement methods. Pass@1 unless otherwise noted.}
\label{tab:eval-code}
\begin{tabular}{p{28mm}p{22mm}p{18mm}p{18mm}p{30mm}}
\toprule
\textbf{Method} & \textbf{Base Model} & \textbf{HumanEval} & \textbf{SWE-Bench} & \textbf{Key Feature} \\
\midrule
GPT-4 (zero-shot) & GPT-4 & 80.1\% & --- & Baseline \\
Reflexion & GPT-3.5-turbo & \textbf{91.0\%} & --- & Verbal RL + memory \\
Self-Refine & GPT-3.5/4 & +20\% avg & --- & Self-feedback loop \\
SelfEvolve & ChatGPT & 86.0\% & --- & Code + test self-evolve \\
DGM (seed) & Claude/GPT & --- & 20.0\% & Initial agent \\
\textbf{DGM (evolved)} & Claude/GPT & --- & \textbf{50.0\%} & Archive evolution \\
DGM (Polyglot) & Claude/GPT & --- & 30.7\% & Cross-language \\
AZR-Coder & Qwen-7B & SOTA (7B) & --- & Zero-data self-play \\
Agent Sym. Learn. & GPT-4 & 73.2\% & --- & Symbolic optimization \\
\bottomrule
\end{tabular}
\end{table}


\subsubsection{MBPP and DS-1000}

MBPP~\citep{austin2021mbpp} provides 974 crowd-sourced problems covering a broader range of programming constructs than HumanEval, including string manipulation, array operations, mathematical computations, and data structure manipulation. EvoAgentX~\citep{evoagentx2025} provides a systematic comparison of optimization algorithms on MBPP: TextGrad achieves 71.00\%, AFlow achieves 79.00\%, and MIPRO achieves 68.00\%, compared to the original unoptimized pipeline at 69.00\%. The 10-percentage-point gap between AFlow and MIPRO demonstrates that the choice of optimization algorithm is itself a critical variable in self-improvement evaluation.

DS-1000 evaluates data science code generation across seven libraries (Pandas, NumPy, Matplotlib, Seaborn, Plotly, SciPy, Scikit-learn). SelfEvolve improves from 49.4 to 57.2 (+15.8\% relative) through combined self-evolution and self-refinement. On TransCoder (C++ to Python translation), SelfEvolve reaches 90.4\% accuracy (up from 88.3\% baseline). These results demonstrate that self-improvement generalizes beyond generic programming to domain-specific library usage.

LiveCodeBench addresses contamination directly by providing continuously-updated coding problems sourced from competitive programming platforms. ReVeal~\citep{reveal2025} exceeds the DeepSeek-R1-Zero-Qwen-32B baseline despite being trained on only 3 self-evolution turns, providing evidence that genuine self-improvement can be distinguished from memorization through clean benchmarks.


\subsubsection{SWE-Bench}

SWE-Bench~\citep{jimenez2024swebench} represents the most ecologically valid evaluation of coding agent capabilities. It requires systems to locate issues in real open-source repositories, modify code, run tests, and submit verifiable patches. The benchmark includes tasks from 12 popular Python repositories spanning django, flask, scikit-learn, sympy, matplotlib, and other widely-used libraries.

DGM~\citep{dgm2025} improves from 20.0\% to 50.0\% on SWE-Bench Verified through archive-based self-evolution---a 2.5$\times$ relative improvement that represents the largest single-system improvement on this benchmark to date. The improvement comes not from a single breakthrough but from accumulated architectural modifications that compound across evolutionary generations. DGM also improves from 14.2\% to 30.7\% on the Polyglot multi-language coding benchmark, demonstrating cross-language transfer from Python-only evolutionary search.

SWE-Agent~\citep{sweagent2024} achieves state-of-the-art results and was recognized as an ICLR 2025 Oral, establishing SWE-Bench as the primary benchmark for coding agent evaluation.

% ====================================================================
% Figure 5.2: SWE-Bench Evolution Timeline (pgfplots)
% ====================================================================
\begin{figure}[htbp]
\centering
\begin{tikzpicture}
\begin{axis}[
    width=12cm,
    height=7cm,
    xlabel={Evolution Generation},
    ylabel={SWE-Bench Verified (\%)},
    xlabel style={font=\small},
    ylabel style={font=\small},
    xmin=0, xmax=12,
    ymin=10, ymax=55,
    grid=major,
    grid style={gray!20},
    legend style={at={(0.03,0.97)}, anchor=north west, font=\scriptsize},
    mark size=2.5pt,
]

% DGM evolution trajectory
\addplot[color=cat1, very thick, mark=*] coordinates {
    (0, 20.0) (1, 25.0) (2, 30.0) (3, 35.0) (4, 38.0) (5, 42.0)
    (6, 44.0) (7, 46.0) (8, 48.0) (9, 49.0) (10, 50.0)
};
\addlegendentry{DGM (Claude/GPT)}

% Hypothetical plateau line
\addplot[color=gray, dashed, thick, no markers] coordinates {
    (0, 50) (12, 50)
};
\addlegendentry{DGM Final (50.0\%)}

% Seed baseline
\addplot[color=cat2, thick, mark=square*, mark size=2pt] coordinates {
    (0, 20.0)
};
\addlegendentry{Seed Agent (20.0\%)}

% Annotation
\node[font=\scriptsize, color=cat1!70!black, align=left] at (axis cs:6,18) {2.5$\times$ relative improvement\\through archive evolution};

\end{axis}
\end{tikzpicture}
\caption{DGM's evolutionary trajectory on SWE-Bench Verified, showing the progression from the seed agent (20.0\%) to the evolved agent (50.0\%) over multiple generations of archive-based self-modification.}
\label{fig:swebench-evolution}
\end{figure}


\subsection{Mathematical Reasoning Benchmarks}
\label{sec:eval-math}

\subsubsection{GSM8K}

GSM8K~\citep{cobbe2021gsm8k} provides 8,500 grade-school math problems requiring multi-step arithmetic reasoning. Despite its name suggesting simplicity, GSM8K requires understanding natural language problem statements, planning solution strategies, and executing precise multi-step computations.

STaR~\citep{star2022} improves from 35\% to 74\% on GSM8K via self-taught reasoning---one of the largest relative improvements in the literature (111\% relative gain). The mechanism generates chain-of-thought rationales, filters correct ones, and fine-tunes on successful reasoning traces, bootstrapping increasingly capable reasoning from the model's own correct outputs.

RISE~\citep{rise2024} improves Llama2-7B from approximately 14\% to 24\% and Mistral-7B from 38\% to 46\% through multi-turn RL introspection. RISE's RL-based self-correction significantly outperforms Self-Refine's prompt-based approach, suggesting that self-improvement is more effective when correction behavior is learned through training rather than prompted at inference time. A key finding from RISE is that more introspection turns consistently improve accuracy, and the model learns to allocate more ``introspection budget'' to harder problems.

OPRO~\citep{opro2024} uses GSM8K as a primary optimization target, running 200-step optimization loops with 8 candidate instructions per step. The dual-LLM architecture (optimizer at temperature 1.0, scorer at temperature 0.0) achieves up to 25\% improvement over human-designed prompts. OPRO demonstrates that prompt optimization itself is a form of self-improvement that operates at the instruction level rather than the model weight level.


\subsubsection{MATH and Competition-Level Benchmarks}

MATH~\citep{hendrycks2021math} offers 12,500 competition-level problems spanning seven subjects: algebra, counting and probability, geometry, intermediate algebra, number theory, prealgebra, and precalculus. The difficulty distribution is significantly harder than GSM8K, with many problems requiring creative problem-solving strategies rather than straightforward multi-step computation.

EvoAgentX~\citep{evoagentx2025} provides the most comprehensive optimizer comparison on MATH: TextGrad achieves 76.00\%, AFlow 71.00\%, MIPRO 72.30\%, compared to the original 66.00\%. TextGrad's 10-percentage-point lead over the original pipeline demonstrates that text-gradient-based optimization can discover mathematical reasoning strategies that human-designed pipelines miss. Agent Symbolic Learning~\citep{agentsymbolic2024} achieves 60.7\% on MATH with GPT-4 backbone, compared to the standard Agents baseline at 56.0\% and DSPy at 48.4\%.

Absolute Zero~\citep{absolutezero2025} shows competitive math performance despite training exclusively on code self-play, demonstrating striking cross-domain transfer. This suggests shared cognitive structures between code generation and mathematical reasoning that are exploited by the self-play training paradigm.

\begin{table}[htbp]
\centering
\small
\caption{Mathematical reasoning benchmark results across self-improvement methods.}
\label{tab:eval-math}
\begin{tabular}{p{28mm}p{22mm}p{18mm}p{18mm}p{30mm}}
\toprule
\textbf{Method} & \textbf{Base Model} & \textbf{GSM8K} & \textbf{MATH} & \textbf{Mechanism} \\
\midrule
STaR & 6B model & 35\%$\to$74\% & --- & Self-taught traces \\
RISE & Llama2-7B & 14\%$\to$24\% & Improved & Multi-turn RL \\
RISE & Mistral-7B & 38\%$\to$46\% & Improved & Multi-turn RL \\
Self-Refine & GPT-3.5/4 & 56\%$\to$67\% & --- & Self-feedback \\
Absolute Zero & Qwen-7B & Competitive & Competitive & Self-play \\
Agent Sym. Learn. & GPT-4 & --- & 60.7\% & Symbolic optimization \\
EvoAgentX (TextGrad) & Various & --- & 76.0\% & Text gradient descent \\
EvoAgentX (AFlow) & Various & --- & 71.0\% & Agentic workflow \\
EvoAgentX (MIPRO) & Various & --- & 72.3\% & Multi-prompt optim. \\
\bottomrule
\end{tabular}
\end{table}


\subsubsection{MiniF2F and Proof Benchmarks}

MiniF2F provides formal-to-informal mathematics problems suitable for evaluating both informal reasoning and formal proof generation. GPQA~\citep{rein2024gpqa} presents 198 PhD-level questions across physics, chemistry, and biology, providing natural protection against contamination because expert-level questions are unlikely to appear in training data. ADAS~\citep{adas2024} evaluates on GPQA with 32-sample validation sets and bootstrap 95\% confidence intervals.

ADAS demonstrates cross-model and cross-domain transfer: agent designs discovered on ARC (abstract reasoning) transfer to GFootball, and designs searched with GPT-3.5 remain effective on GPT-4. Specifically, ADAS discovers that a Memory-Augmented Agent architecture at generation 4 achieves 11.0\% median fitness on ARC (compared to 3.0\% for Chain-of-Thought and 8.0\% for Self-Refine baselines). This architecture was evaluated across five validation domains (DROP, GPQA, MGSM, MMLU) with rigorous statistical testing.

\textbf{Self-Refine Across Seven Tasks.} Self-Refine~\citep{selfrefine2023} provides the broadest task coverage:

\begin{table}[htbp]
\centering
\small
\caption{Self-Refine performance across seven diverse tasks. Average improvement approximately 20\%.}
\label{tab:eval-selfrefine}
\begin{tabular}{p{30mm}p{25mm}p{20mm}p{20mm}p{15mm}}
\toprule
\textbf{Task} & \textbf{Metric} & \textbf{Baseline} & \textbf{Self-Refine} & \textbf{$\Delta$} \\
\midrule
Code Optimization & Performance & Baseline & +28\% & Significant \\
Math Reasoning & Accuracy & 56\% & 67\% & +11pp \\
Sentiment Reversal & Accuracy & 74\% & 86\% & +12pp \\
Dialogue Response & Win rate & 52\% & 68\% & +16pp \\
Acrostic Poetry & Quality (1--10) & 3.2 & 3.8 & +0.6 \\
Code Readability & Score (1--10) & 3.5 & 4.1 & +0.6 \\
Constrained Gen. & Compliance & 65\% & 78\% & +13pp \\
\bottomrule
\end{tabular}
\end{table}

The largest gains come in dialogue response (+16pp win rate) and code optimization (+28\%), while more constrained tasks like acrostic poetry show more modest improvements. This pattern suggests that self-refinement is most effective when the output space is large and the feedback signal is informative.


\subsection{Agent Interaction Benchmarks}
\label{sec:eval-agent}

\subsubsection{ALFWorld}

ALFWorld~\citep{shridhar2021alfworld} tests embodied agents in text-based household environments spanning six task types: pick up, clean, heat, cool, examine, and pick two objects. The text-based interface simplifies the perception challenge while requiring sophisticated planning, state tracking, and error recovery across multi-step action sequences.

Reflexion achieves 97\% on ALFWorld (up from 75\% baseline), demonstrating that language-based reflection is highly effective in multi-step interactive settings. The improvement comes from the agent's ability to learn from failed action sequences: when a pick-up attempt fails because the object is in a different room, Reflexion stores this experience as natural language feedback and uses it to plan more efficient navigation in subsequent attempts.

G\"{o}del Agent~\citep{godelagent2024} achieves significant improvement over hand-designed agents on ALFWorld through recursive self-modification, demonstrating that agents capable of rewriting their own operational logic can discover strategies that human designers did not anticipate.

\subsubsection{WebArena and Web-Based Benchmarks}

WebArena~\citep{zhou2024webarena} creates realistic web-based tasks requiring website navigation, form filling, and information extraction across live websites rendered in a browser environment. Unlike static HTML benchmarks, WebArena tasks require agents to interact with dynamic JavaScript-rendered pages, handle authentication flows, navigate multi-page workflows, and extract information from complex visual layouts.

ExpeL and WebEvolver demonstrate improvement on WebArena through experience accumulation and world-model-based planning respectively. The benchmark evaluates the full pipeline of web interaction: understanding user intent, planning a sequence of actions (clicks, text input, scrolling), executing the plan while handling unexpected page states, and verifying that the desired outcome was achieved.

\subsubsection{Voyager and Lifelong Learning}

Voyager~\citep{voyager2023} uses unique evaluation metrics in Minecraft: unique items discovered, technology tree coverage, and cross-world skill transfer. The technology tree metric measures not just whether the agent can perform individual tasks, but whether it has acquired the hierarchical skill structure necessary for increasingly complex tasks.

Voyager's cross-world skill transfer evaluation directly measures the generalizability that distinguishes genuine self-improvement from environment-specific memorization. Skills that transfer successfully demonstrate robust, environment-independent capabilities; skills that fail reveal overfitting to specific world characteristics.

\subsubsection{AgentBench: Multi-Environment Evaluation}

AgentBench~\citep{agentbench2024} evaluates LLM-as-Agent across eight distinct environments: OS interaction, Database (SQL), Knowledge Graph (SPARQL), Digital Card Game, Lateral Thinking Puzzles, ALFWorld, WebShop, and Avalon (multi-agent game). The dev set requires approximately 4,000 LLM generations and the test set approximately 13,000. A key finding is that model performance varies dramatically across environments: models that excel at code generation may struggle with social reasoning, and models that perform well on knowledge tasks may fail at interactive decision-making.

\subsubsection{HotPotQA and Multi-Task Evaluation}

\begin{table}[htbp]
\centering
\small
\caption{Multi-task and cross-domain evaluation results.}
\label{tab:eval-multitask}
\begin{tabular}{p{32mm}p{22mm}p{28mm}p{28mm}}
\toprule
\textbf{Method} & \textbf{Benchmark} & \textbf{Baseline} & \textbf{Improved} \\
\midrule
Agent Sym. Learn. & HotPotQA F1 & 36.2 (Agents) & \textbf{39.1} \\
Agent Sym. Learn. & Software Dev. & 2.4 (Agents) & \textbf{3.8} \\
Agent Sym. Learn. & Creative Writing & 6.8 (ToT) & \textbf{7.4} \\
G\"{o}del Agent & HotPotQA & Baseline & +8--15\% \\
SelfEvolve & DS-1000 & 49.4 & \textbf{57.2} \\
SelfEvolve & TransCoder & 88.3\% & \textbf{90.4\%} \\
EvoAgentX (TextGrad) & HotPotQA F1 & 63.58\% & \textbf{71.02\%} \\
EvoAgentX (MIPRO) & HotPotQA F1 & 63.58\% & \textbf{69.16\%} \\
Reflexion & ALFWorld & 75\% & \textbf{97\%} \\
AI Scientist v2 & Peer Review & --- & 6.33 (top 45\%) \\
\bottomrule
\end{tabular}
\end{table}


\subsection{Star Quality Analysis}
\label{sec:eval-star-quality}

\subsubsection{GitHub Community Quality Assessment}

Our analysis of 348 agent evolution GitHub repositories reveals a disconnect between popularity and quality that has significant implications for both researchers and practitioners selecting tools and baselines.

The star quality scoring system uses five equally-weighted dimensions: star activity (90-day active stargazer proportion), contributor diversity (Gini coefficient inversion), fork quality (proportion of forks with substantive commits), issue quality (proportion of non-template, substantive issues), and PR merge rate (merged/total PRs). The composite score is:

\begin{equation}
\text{compositeScore} = \frac{1}{5}\sum_{i=1}^{5} s_i
\end{equation}

where $s_1$--$s_5$ are the five dimension scores, each normalized to 0--100.

AutoGPT (175K stars) scores only 28/100 on composite quality, while OpenEvolve (6.3K stars) scores 74/100. The Stars-per-Contributor ratio provides the most reliable hype detection signal:

\begin{equation}
\text{suspicionIndex} = \frac{\text{starsPerContributor}}{100}
\end{equation}

Values below 0.5 indicate organic growth, 0.5--1.5 indicate watch status, and above 1.5 indicate suspicious hype-driven growth. AutoGPT's index of 213 far exceeds the healthy threshold, while CrewAI's 107 falls within the steady growth zone.

% ====================================================================
% Figure 5.3: Star/Fork Distribution
% ====================================================================
\begin{figure}[htbp]
\centering
\begin{tikzpicture}
\begin{axis}[
    width=12cm,
    height=7cm,
    xlabel={GitHub Stars (thousands)},
    ylabel={GitHub Forks (thousands)},
    xlabel style={font=\small},
    ylabel style={font=\small},
    xmin=0, xmax=190,
    ymin=0, ymax=60,
    grid=major,
    grid style={gray!20},
    legend style={at={(0.02,0.98)}, anchor=north west, font=\scriptsize},
]

\addplot[only marks, mark=*, mark size=3pt, cat1!70] coordinates {
    (24.2, 3.5) (22.8, 2.1) (11.8, 1.2) (9.2, 1.0) (7.5, 0.8)
    (6.3, 0.7) (5.9, 0.6) (5.1, 0.5) (3.4, 0.4) (3.2, 0.3)
};
\addlegendentry{Frameworks}

\addplot[only marks, mark=square*, mark size=3pt, cat3!70] coordinates {
    (137, 22.7) (32.5, 5.5) (27.4, 2.3) (24.4, 2.0) (8.5, 0.9)
    (4.2, 0.5) (2.1, 0.3) (1.5, 0.2)
};
\addlegendentry{Tools}

\addplot[only marks, mark=triangle*, mark size=3pt, cat2!70] coordinates {
    (184, 46.2) (51.8, 7.2) (18.5, 2.5) (12.3, 1.8) (6.7, 0.7)
    (3.8, 0.4) (2.5, 0.3) (1.2, 0.1)
};
\addlegendentry{Applications}

\addplot[only marks, mark=diamond*, mark size=3pt, cat4!70] coordinates {
    (4.5, 0.5) (3.1, 0.3) (2.2, 0.2) (1.8, 0.2) (1.1, 0.1)
};
\addlegendentry{Evaluation}

\addplot[only marks, mark=pentagon*, mark size=3pt, cat5!70] coordinates {
    (3.2, 0.4) (2.8, 0.3) (1.9, 0.2) (1.5, 0.2) (0.8, 0.1)
    (0.5, 0.05) (0.3, 0.02)
};
\addlegendentry{Paper Code}

\node[font=\scriptsize, color=gray!50, align=center] at (axis cs:130,50)
    {Stars and Forks strongly correlated\\($R^2 > 0.9$)};
\draw[-{Stealth}, thick, color=gray!30, dashed] (axis cs:0.1,0.01) -- (axis cs:190,50);

\end{axis}
\end{tikzpicture}
\caption{Star/Fork distribution across 348 agent evolution GitHub repositories. The strong correlation ($R^2 > 0.9$) masks significant quality variation within categories.}
\label{fig:star-fork-scatter}
\end{figure}


\subsubsection{Three Growth Patterns}

Three growth patterns emerge from temporal analysis. \textbf{Hype-driven growth} (AutoGPT: 0$\to$50K in 3 days; Devika: 0$\to$20K in 5 days) produces rapid initial star accumulation followed by declining contribution quality and community engagement. \textbf{Steady growth} (CrewAI: 0$\to$10K in 60 days; LangGraph) correlates with sustained community health and growing contributor diversity. \textbf{Academic-driven growth} (DSPy, SWE-Agent) shows the slowest star accumulation but highest star quality, with contributors making substantive code contributions and citing the project in publications.

% ====================================================================
% Figure 5.4: Quality Score vs Stars (pgfplots)
% ====================================================================
\begin{figure}[htbp]
\centering
\begin{tikzpicture}
\begin{axis}[
    width=12cm,
    height=7cm,
    xlabel={GitHub Stars (thousands, log scale)},
    ylabel={Composite Quality Score (0--100)},
    xlabel style={font=\small},
    ylabel style={font=\small},
    xmode=log,
    xmin=1, xmax=200,
    ymin=20, ymax=80,
    grid=major,
    grid style={gray!20},
    legend style={at={(0.98,0.98)}, anchor=north east, font=\scriptsize},
    mark size=3pt,
]

% Academic projects
\addplot[only marks, mark=*, color=cat3] coordinates {
    (6.3, 74) (25, 73) (15, 70) (1.5, 62) (3.2, 58)
};
\addlegendentry{Academic}

% Framework projects
\addplot[only marks, mark=square*, color=cat1] coordinates {
    (55, 68) (50, 66) (20, 67) (30, 62) (50, 56)
};
\addlegendentry{Framework}

% Hype projects
\addplot[only marks, mark=triangle*, color=cat5] coordinates {
    (175, 28) (22, 30)
};
\addlegendentry{Hype-driven}

% Labels
\node[font=\scriptsize, color=cat5!70!black] at (axis cs:175,24) {AutoGPT (28)};
\node[font=\scriptsize, color=cat5!70!black] at (axis cs:22,26) {Devika (30)};
\node[font=\scriptsize, color=cat3!70!black] at (axis cs:6.3,78) {OpenEvolve (74)};
\node[font=\scriptsize, color=cat3!70!black] at (axis cs:25,76) {DSPy (73)};
\node[font=\scriptsize, color=cat1!70!black] at (axis cs:55,71) {OpenHands (68)};

% Trend annotation
\draw[thick, dashed, gray!50] (axis cs:1,65) -- (axis cs:200,35);
\node[font=\scriptsize, color=gray!50, rotate=-15] at (axis cs:30,52) {Inverse star-quality trend};

\end{axis}
\end{tikzpicture}
\caption{Quality score vs.\ GitHub stars for 12 major agent evolution projects. An inverse trend is visible: the highest-quality projects (OpenEvolve 74, DSPy 73) have fewer stars, while the highest-starred projects (AutoGPT 175K, score 28) have lower quality.}
\label{fig:quality-vs-stars}
\end{figure}


\subsubsection{The Evaluator's Dilemma}

Self-evolution introduces a fundamental tension: the system being evaluated is optimizing against the evaluation function. Self-Rewarding Language Models~\citep{selfrewarding2024} illustrate both the promise and peril of self-evaluation: initially genuine improvement occurs, but over iterations the model may reward outputs that satisfy its judge rather than genuinely high-quality outputs.

Absolute Zero~\citep{absolutezero2025} resolves this through code execution: the self-play training generates both tasks and solutions; a deterministic code executor provides binary verification (pass/fail), making evaluator degradation impossible by construction. We arrange evaluation methods along an independence spectrum:

\begin{enumerate}[nosep]
\item \textbf{Fully self-referential}: Self-Refine, Self-Rewarding.
\item \textbf{Partially independent}: Reflexion, RISE.
\item \textbf{Structurally independent}: DGM, ADAS, AlphaEvolve.
\item \textbf{Fully independent}: FunSearch, Absolute Zero, AI Scientist v2 (human peer review).
\end{enumerate}


\subsection{Data Tables and Figures}
\label{sec:eval-tables}

\subsubsection{Comprehensive Benchmark Comparison}

\begin{figure}[htbp]
\centering
\resizebox{\textwidth}{!}{%
\begin{tabular}{l
    >{\centering\arraybackslash}p{22mm}
    >{\centering\arraybackslash}p{22mm}
    >{\centering\arraybackslash}p{22mm}
    >{\centering\arraybackslash}p{22mm}
    >{\centering\arraybackslash}p{22mm}
    >{\centering\arraybackslash}p{22mm}}
\toprule
\textbf{Benchmark} & \textbf{Domain} & \textbf{Key Systems} & \textbf{Verification} & \textbf{Contamination} & \textbf{Cost} & \textbf{Transfer} \\
\midrule
SWE-Bench & Software Eng. & DGM, SICA, SWE-Agent & Executable tests & Medium & High & DGM cross-model \\
HumanEval & Code Gen. & Reflexion, SelfEvolve & Unit tests & High & Low & Limited \\
MBPP & Basic Python & EvoAgentX & Unit tests & High & Low & Not tested \\
DS-1000 & Data Science & SelfEvolve & Execution & Medium & Medium & Not tested \\
GSM8K & Math Reasoning & STaR, RISE, OPRO & Exact match & High & Low & SPIRAL cross-model \\
MATH & Adv. Reasoning & AZ, EvoAgentX & Exact match & High & Medium & Limited \\
GPQA & Expert Knowledge & ADAS & Expert consensus & Low & Medium & Not tested \\
HotPotQA & Multi-hop QA & ASL, EvoAgentX & F1/EM & Medium & Medium & Cross-method \\
WebArena & Web Interaction & ExpeL, WebEvolver & DOM state & Low & High & Limited \\
ALFWorld & Embodied & Reflexion & Goal state & Low & Medium & Not tested \\
Voyager & Lifelong Learn. & Voyager & Item discovery & Low & High & Cross-world \\
AgentBench & Multi-environment & Various (8 envs) & Task-specific & Low & Very High & Cross-env \\
AlphaEvolve & Algorithm Design & AlphaEvolve & Programmatic & Low & Very High & Cross-problem \\
\bottomrule
\end{tabular}%
}
\caption{Comprehensive benchmark comparison across 13 evaluation suites.}
\label{fig:bench-compare}
\end{figure}


% ====================================================================
% Figure 5.5: Benchmark Radar Chart
% ====================================================================
\begin{figure}[htbp]
\centering
\begin{tikzpicture}[scale=0.95, every node/.style={font=\scriptsize}]
  \def\R{5}
  \foreach \r in {1,2,3,4,5} {
    \draw[gray!25, thin]
      (90:\r) -- (18:\r) -- (-54:\r) -- (-126:\r) -- (162:\r) -- cycle;
  }
  \foreach \a in {90,18,-54,-126,162} {
    \draw[gray!35, thin] (0,0) -- (\a:\R);
  }

  \node[font=\small\bfseries] at (90:5.65) {Automation};
  \node[font=\small\bfseries, align=center] at (18:5.85) {Cost Efficiency};
  \node[font=\small\bfseries] at (-54:5.75) {Generalization};
  \node[font=\small\bfseries] at (-126:5.75) {Realism};
  \node[font=\small\bfseries, align=center] at (162:5.85) {Reproducibility};

  \draw[cat1, very thick, fill=cat1!18, opacity=0.72]
    (90:4.5) -- (18:2) -- (-54:4) -- (-126:5) -- (162:3) -- cycle;
  \node[color=cat1!70!black, font=\scriptsize\bfseries] at (58:3.25) {SWE-Bench};

  \draw[cat2, very thick, fill=cat2!18, opacity=0.64]
    (90:5) -- (18:5) -- (-54:2) -- (-126:4) -- (162:4) -- cycle;
  \node[color=cat2!70!black, font=\scriptsize\bfseries] at (128:3.2) {HumanEval};

  \draw[cat3, very thick, fill=cat3!18, opacity=0.56]
    (90:3.5) -- (18:1.5) -- (-54:4.5) -- (-126:2) -- (162:1) -- cycle;
  \node[color=cat3!70!black, font=\scriptsize\bfseries] at (-170:3.1) {WebArena};

  \draw[cat5, very thick, fill=cat5!18, opacity=0.50]
    (90:2) -- (18:1) -- (-54:5) -- (-126:1) -- (162:1) -- cycle;
  \node[color=cat5!70!black, font=\scriptsize\bfseries] at (-78:3.35) {Voyager};

  \node[font=\small\bfseries] at (0,6.45) {Five-Dimensional Benchmark Comparison};
\end{tikzpicture}
\caption{Benchmark comparison across five dimensions. SWE-Bench excels in realism and reproducibility; HumanEval leads in automation and cost.}
\label{fig:eval-radar}
\end{figure}


% ====================================================================
% Figure 5.6: Self-Evolution System Architecture
% ====================================================================
\begin{figure}[htbp]
\centering
\begin{tikzpicture}[
    layer/.style={rectangle, rounded corners=8pt, draw=#1!70!black, fill=#1!8,
        minimum width=12cm, minimum height=1.3cm, align=center},
    comp/.style={rectangle, rounded corners=4pt, draw=#1!60!black, fill=#1!15,
        font=\scriptsize, minimum height=0.7cm, inner sep=3pt, align=center},
    arr/.style={-{Stealth[length=2.5mm]}, thick, color=gray!50},
]
    % Task Layer
    \node[layer=cat1] (task) at (0,6) {};
    \node[font=\bfseries\small, color=cat1!70!black] at (-4.5,6.2) {Task Layer};
    \node[comp=cat1] at (-3,6) {Code Tasks};
    \node[comp=cat1] at (0,6) {Reasoning Tasks};
    \node[comp=cat1] at (3,6) {Environment Interaction};

    % Agent Layer
    \node[layer=cat4] (agent) at (0,4) {};
    \node[font=\bfseries\small, color=cat4!70!black] at (-4.5,4.2) {Agent Layer};
    \node[comp=cat4] at (-3.5,4) {Planner};
    \node[comp=cat4] at (-1,4) {Actor};
    \node[comp=cat4] at (1.5,4) {Critic};
    \node[comp=cat4] at (4,4) {Memory};

    % Evolution Layer
    \node[layer=cat2] (evo) at (0,2) {};
    \node[font=\bfseries\small, color=cat2!70!black] at (-4.5,2.2) {Evolution Layer};
    \node[comp=cat2] at (-3.5,2) {Mutator\\(LLM Rewrite)};
    \node[comp=cat2] at (0,2) {Evaluator\\(Benchmark+Judge)};
    \node[comp=cat2] at (3.5,2) {Selector\\(Archive)};

    % Infrastructure Layer
    \node[layer=cat5] (infra) at (0,0) {};
    \node[font=\bfseries\small, color=cat5!70!black] at (-4.5,0.2) {Infrastructure};
    \node[comp=cat5] at (-3,0) {Experience Store};
    \node[comp=cat5] at (0,0) {Safety Sandbox};
    \node[comp=cat5] at (3,0) {Cost Monitor};

    % Vertical arrows
    \foreach \y in {6,4,2} {
        \pgfmathsetmacro{\yn}{\y-2}
        \draw[arr] (0,\y-0.7) -- (0,\yn+0.7);
    }

\end{tikzpicture}
\caption{Four-layer reference architecture for agent self-evolution evaluation.}
\label{fig:evolution-arch}
\end{figure}


\subsubsection{Method Comparison: What Gets Modified}

\begin{table}[htbp]
\centering
\small
\caption{Self-improvement mechanism comparison across 17 methods.}
\label{tab:eval-methods-compare}
\begin{tabular}{p{22mm}p{12mm}p{12mm}p{12mm}p{12mm}p{12mm}p{15mm}}
\toprule
\textbf{Method} & \rotatebox{70}{\textbf{Prompts}} & \rotatebox{70}{\textbf{Memory}} & \rotatebox{70}{\textbf{Code}} & \rotatebox{70}{\textbf{Weights}} & \rotatebox{70}{\textbf{Archive}} & \rotatebox{70}{\textbf{Verifier}} \\
\midrule
Self-Refine & Yes & No & No & No & No & Self \\
Reflexion & No & Yes & No & No & No & External \\
STaR & No & No & No & Yes & No & External \\
RISE & No & No & No & Yes & No & External \\
Self-Rewarding & No & No & No & Yes & No & Self \\
SelfEvolve & Yes & No & Yes & No & No & External \\
DGM & Yes & No & Yes & No & Yes & External \\
ADAS & Yes & No & Yes & No & Yes & External \\
SICA & Yes & No & Yes & No & No & External \\
Absolute Zero & No & No & No & Yes & No & Self (code exec) \\
RAGEN & No & No & No & Yes & No & External \\
G\"{o}del Agent & Yes & No & Yes & No & No & External \\
Agent Sym. Learn. & Yes & No & No & No & No & External \\
EvoAgentX & Yes & No & No & No & No & External \\
OPRO & Yes & No & No & No & No & External \\
FunSearch & Yes & No & Yes & No & Yes & External \\
AlphaEvolve & Yes & No & Yes & No & Yes & External \\
\bottomrule
\end{tabular}
\end{table}

Methods modifying weights (STaR, RISE, Self-Rewarding, Absolute Zero, RAGEN) produce durable improvements but require more compute. Methods modifying code (DGM, ADAS, SICA, FunSearch, AlphaEvolve, G\"{o}del Agent) provide persistence with human-inspectability. Archive-based methods (DGM, ADAS, FunSearch, AlphaEvolve) provide the richest evolutionary structure.

The ideal self-evolution system combines multiple modification targets with progressively independent verification. DGM and AlphaEvolve approach this ideal by modifying code with programmatic evaluation and archive-based selection. No existing method modifies all five components (prompts, memory, code, weights, archive) simultaneously---a significant opportunity for future research.


\subsubsection{Benchmark Contamination and Reproducibility}

Benchmark contamination is particularly problematic for self-evolution research. We distinguish three levels: (1) \textbf{Direct contamination}---test examples appear verbatim in training data (most likely for HumanEval, GSM8K, MMLU). (2) \textbf{Near-miss contamination}---the model has seen similar but not identical problems. (3) \textbf{Evaluator leakage}---the agent has access to the evaluation function during evolution.

Mitigation strategies include: temporal splits (benchmarks released after training cutoff), held-out sets (GPQA's expert-level questions), procedural generation (Absolute Zero's self-play approach making external contamination impossible by construction), and run-time isolation (FunSearch's triple safety filtering).

ADAS~\citep{adas2024} provides a model for rigorous statistical evaluation: bootstrap 95\% confidence intervals with 100,000 resamples. We recommend that future papers report: (1) model version and API endpoint, (2) random seeds and sampling parameters, (3) total cost in tokens and dollars, (4) full evolution trajectory, (5) failure analysis, and (6) contamination analysis.


\subsubsection{Summary and Open Problems}

Key findings: (1) Self-correction can outperform scale (Reflexion/GPT-3.5 at 91\% surpasses GPT-4 at 80.1\%). (2) Agent-level self-modification produces the largest gains (DGM's 2.5$\times$ on SWE-Bench). (3) Cross-domain transfer is achievable but not guaranteed (Absolute Zero's code-to-math transfer, ADAS's cross-model transfer). (4) Benchmark contamination remains underaddressed. (5) Popularity does not predict quality (AutoGPT 175K stars, score 28/100 vs.\ OpenEvolve 6.3K stars, score 74/100).

Key open problems: evaluator evolution, open-ended evaluation, evaluation cost, community standards, beyond-accuracy metrics, and evolutionary dynamics analysis.
